% TOP_PROBLEM: {"id": 30007152, "title": "Bounded holomorphic functions under one-sided negative curvature", "category_id": 6, "category": {"id": 6, "name": "geometry", "display_name": "Geometry", "description": "Euclidean and non-Euclidean geometry, geometric structures.", "slug": "geometry", "order_index": 6, "created_at": "2026-07-31T15:26:25.671Z"}, "status": "open"}
% FIRST_POSTED: 2026-10-06
% ATTEMPT_STATUS: solved
% OPENAI_CLAIM: {"schema_version":1,"source_repo":"https://github.com/openai/math","source_commit":"adc7f1241b42e322a6451854ab7e4b4c146bf78a","release_date":"2026-10-06","match":"full","resolution":"disproved","independently_reviewed":false,"second_pass":"agreed","families":[{"family":"359","title":"Negative Kähler curvature without bounded holomorphic coordinates","subject":"Differential geometry","manuscripts":[{"title":"One-sided negative sectional curvature and the holomorphic Liouville property","pdf_path":"preprints/One-sided-negative-sectional-curvature-and-the-holomorphic-Liouville-property-September-25-2026/paper.pdf","readme_path":"preprints/One-sided-negative-sectional-curvature-and-the-holomorphic-Liouville-property-September-25-2026/README.md","date":"2026-09-25","theorem_ref":"thm:main (01-introduction.tex:21)","statement_sources":[{"path":"preprints/One-sided-negative-sectional-curvature-and-the-holomorphic-Liouville-property-September-25-2026/build/01-introduction.tex","line":21,"label":"thm:main"}]}],"lean":{"doc_path":"lean/docs/359.md","comparators":[],"declarations":[],"covers_main_theorem":false},"reasoning_trace":null}]}
\documentclass[11pt]{article}
\usepackage[margin=25mm]{geometry}
\usepackage{fontspec}
\setmainfont{Latin Modern Roman}
\usepackage{amsmath,amssymb,mathtools}
\usepackage{unicode-math}
\setmathfont{Latin Modern Math}
\usepackage{enumitem}
\usepackage{xurl,hyperref,bookmark}
\hypersetup{colorlinks=true,urlcolor=blue}
\setcounter{secnumdepth}{0}
\setlength{\parindent}{0pt}
\setlength{\parskip}{5pt}
\setlength{\emergencystretch}{3em}
\begin{document}
\section*{Bounded holomorphic functions under one-sided negative curvature}

\subsection{OpenAI's claim}
\textit{One-sided negative Kähler curvature and Liouville}: The paper claims a contractible complex domain in a fixed finite dimension with a complete Kähler metric of sectional curvature at most minus one and only constant bounded holomorphic functions. Curvature is allowed to be unbounded in the negative direction.

\subsection{Mathematical statement of the claimed result}
\paragraph{Mathematical claim 1.} For some finite complex dimension $m\ge2$, the claim supplies a domain $\mathcal T\subset\mathbb C^m$ carrying a K\"ahler metric which is complete. Moreover $\mathcal T$ is diffeomorphic to $\mathbb R^{2m}$ and satisfies \[\sec_g\le-1,\qquad H^\infty(\mathcal T)=\mathbb C.\] Thus every bounded holomorphic scalar function is constant. Sectional curvature is unbounded below; this is a one-sided negative-curvature example rather than two-sided pinching.

\subsection{Outline of the argument}
\textit{One-sided negative Kähler curvature and Liouville}: A logarithmic tube potential transfers a complete nonpositively curved base metric to a strictly negative tube metric. Homogeneous polynomial systems locate large holomorphic disks above selected finite spatial sets. Disk estimates force derivatives of every bounded holomorphic function to vanish, while the metric construction controls completeness without adding a lower sectional bound.

\subsection{Scope relative to this problem}
\textit{One-sided negative Kähler curvature and Liouville}: The curvature assumption is the one-sided upper bound, so unbounded negative curvature is allowed. The announced Liouville property rules out every nonconstant bounded holomorphic function on the example; it is stronger than merely lacking a bounded coordinate system.

The claimed counterexample lies within the question’s hypotheses and contradicts its universal conclusion.

\subsection{Formal verification}
Family 359 describes Lean coverage for the two-sided negatively pinched threefold obstruction; the one-sided Liouville construction is separate.

\textit{One-sided negative Kähler curvature and Liouville}: No manuscript-specific Comparator is listed in the pinned inventory.

The documented selection does not cover the entire manuscript assertion relevant here.

This site has not run the Lean development. These links describe OpenAI’s selected formal statements and should be read with their stated scope.

\subsection{Status caveat}
OpenAI reports that an internal model produced the manuscripts. Their arguments have not been independently verified as proofs by this site. Any portion outside the documented formal scope may contain mistakes, and the catalogue label reports the announced mathematical claim.

\subsection{Sources}
\begin{itemize}
\item OpenAI, \emph{One-sided negative sectional curvature and the holomorphic Liouville property}, 2026-09-25; citation key \texttt{OAI:One-sided-negative-sectional-curvature-and-the-holomorphic-Liouville-property-September-25-2026}. \href{https://github.com/openai/math/blob/adc7f1241b42e322a6451854ab7e4b4c146bf78a/preprints/One-sided-negative-sectional-curvature-and-the-holomorphic-Liouville-property-September-25-2026/paper.pdf}{PDF (GitHub)}; \href{https://github.com/openai/math/raw/adc7f1241b42e322a6451854ab7e4b4c146bf78a/preprints/One-sided-negative-sectional-curvature-and-the-holomorphic-Liouville-property-September-25-2026/paper.pdf}{PDF (direct)}; \href{https://github.com/openai/math/blob/adc7f1241b42e322a6451854ab7e4b4c146bf78a/preprints/One-sided-negative-sectional-curvature-and-the-holomorphic-Liouville-property-September-25-2026/README.md}{Manuscript page and citation}; \href{https://github.com/openai/math/blob/main/preprints/One-sided-negative-sectional-curvature-and-the-holomorphic-Liouville-property-September-25-2026/paper.pdf}{Latest version}; Statement source: \href{https://github.com/openai/math/blob/adc7f1241b42e322a6451854ab7e4b4c146bf78a/preprints/One-sided-negative-sectional-curvature-and-the-holomorphic-Liouville-property-September-25-2026/build/01-introduction.tex\#L21}{thm:main}.
\item Family 359: \href{https://github.com/openai/math/blob/adc7f1241b42e322a6451854ab7e4b4c146bf78a/lean/docs/359.md}{Lean scope}.
\item \href{https://github.com/openai/math/blob/adc7f1241b42e322a6451854ab7e4b4c146bf78a/CONTENTS.md}{OpenAI catalogue entry}.
\end{itemize}
\end{document}
